\DocumentMetadata{
  pdfversion=2.0,
  pdfstandard=ua-2,
  lang=en-US,
  tagging=on,
  tagging-setup={math/setup=mathml-SE,tabsorder=structure}
}
\documentclass[11pt,letterpaper]{article}
\usepackage[margin=0.68in,includefoot]{geometry}
\usepackage{newcomputermodern}
\usepackage{amsmath,amscd,mathtools}
\usepackage{tikz,adjustbox}
\providecommand{\square}{\mdlgwhtsquare}
\usepackage[hidelinks]{hyperref}
\hypersetup{
  pdftitle={Analysis PhD Exam, January 1993},
  pdfauthor={University of Florida Department of Mathematics},
  pdfsubject={Graduate mathematics examination},
  pdfdisplaydoctitle=true,
  bookmarksopen=true
}
\setcounter{secnumdepth}{0}
\setlength{\parindent}{0pt}
\setlength{\parskip}{0.28em}
\setlength{\emergencystretch}{3em}
\setlength{\leftmargini}{2.15em}
\setlength{\leftmarginii}{2.65em}
\setlength{\leftmarginiii}{3em}
\renewcommand{\labelenumi}{\textbf{\theenumi.}}
\pagestyle{empty}
\raggedbottom
\allowdisplaybreaks
\definecolor{ProblemConcern}{HTML}{7A1F1F}
\newenvironment{problemconcern}{%
  \par\tagstructbegin{tag=Aside}\begingroup\color{ProblemConcern}
}{%
  \par\endgroup\tagstructend\nopagebreak[4]\vspace{0.12em}
}
\newenvironment{examproblems}[1]{%
  \begin{enumerate}%
  \setcounter{enumi}{#1}%
  \setlength{\topsep}{0.35em}%
  \setlength{\partopsep}{0pt}%
  \setlength{\itemsep}{0.48em}%
  \setlength{\parsep}{0.18em}%
}{\end{enumerate}}
\newenvironment{examsubparts}{%
  \begin{enumerate}%
  \setlength{\topsep}{0.18em}%
  \setlength{\partopsep}{0pt}%
  \setlength{\itemsep}{0.16em}%
  \setlength{\parsep}{0.1em}%
}{\end{enumerate}}
\newenvironment{examinstructions}{%
  \par\begingroup\itshape
}{%
  \par\endgroup\vspace{0.18em}
}
\makeatletter
\renewcommand\section{\@startsection{section}{1}{0pt}%
  {-1.25ex plus -.25ex minus -.1ex}%
  {0.55ex plus .1ex}%
  {\normalfont\large\bfseries}}
\renewcommand{\@maketitle}{%
  \newpage\null\vskip -1em%
  {\footnotesize\bfseries
    \href{https://www.ufl.edu/}{University of Florida}, \href{https://math.ufl.edu/}{Department of Mathematics}\par}%
  \vskip -0.5em%
  \tagstructbegin{tag=Title}%
    {\LARGE\bfseries\href{https://gma.math.ufl.edu/exams/analysis-phd/}{Analysis PhD Exam}, January 1993\par}%
  \tagstructend%
  \vskip 0.25em%
  \tagmcbegin{artifact}%
    \hrule height 1pt%
  \tagmcend%
  \vskip 1em%
}
\makeatother
\title{Analysis PhD Exam, January 1993}
\author{}
\date{}
\begin{document}
\maketitle
\thispagestyle{empty}
\begin{examinstructions}
Be sure to present all work in a neat and logical fashion in order to receive credit.
\end{examinstructions}
\begin{examproblems}{0}
\item
State and prove the Radon–Nikodym theorem.
\item
Let \((S,\Sigma,\mu)\) be a positive finite measure space. Prove that \(\left(L^1(S,\Sigma,\mu)\right)^*=L^\infty(S,\Sigma,\mu)\).
\item
Give the construction of the product measure \(\mu\times\nu\) on \(\mathcal S\otimes\mathcal W\), where \((X,\mathcal S,\mu)\) and \((Y,\mathcal W,\nu)\) are finite non negative measure spaces. (Be sure to prove that \(\mu\times\nu\) is countably additive.)
\item
Let \((S,\Sigma,\mu)\) be a positive measure space with \(\mu(S)=1\). Let \(\mathcal A\) and \(\mathcal B\) be a sub~\(\sigma\)-algebras of~\(\Sigma\). Suppose \(\mu(A\cap B)=\mu(A)\mu(B)\) for every \(A\in\mathcal A\) and \(B\in\mathcal B\). Show that \(\int fg\,d\mu=(\int f\,d\mu)(\int g\,d\mu)\), for \(f\geq 0\), \(g\geq 0\), where~\(f\) is \(\mathcal A\)-measurable and~\(g\) is \(\mathcal B\)-measurable.
\item
Let \((S,\Sigma)\) be a measurable space, and suppose~\(\mu\) and~\(\nu\) are positive finite measures on~\(\Sigma\). Suppose \(\mu\ll\nu\) and \(\nu\ll\mu\). Show that 
\[
\frac{d\mu}{d\nu}=\frac{1}{\dfrac{d\nu}{d\mu}}\quad\text{a.e. }\nu.
\]
\item
\begin{problemconcern}
\textbf{Warning: Suspected error.} The underlying space in the opening tuple appears to be printed as \(S_\infty\), but the integral and the later pointwise-convergence statement refer to~\(S\). The intended symbol is uncertain, so the apparent source reading has been retained.
\end{problemconcern}
Let \((S_\infty,\Sigma,\mu)\) be a finite measure space. Let \((f_n)_{n=1}^\infty\) be a sequence of integrable function such that 
\[
\sum_{n=1}^\infty\int_S |f_n|\,d\mu<\infty.
\]
 Does it follow that \(\sum f_n\) converges pointwise on~\(S\) a.e. \(\mu\)? Show all details.
\item
Let~\(f\) be a Lebesgue integrable function on \(\mathbb R\). Let \(\epsilon>0\). Is it possible to find a bounded interval~\(I\) such that whenever~\(E\) is a measurable set such that \(E\cap I=\varnothing\), then 
\[
\left|\int_E f\,dm\right|<\epsilon?
\]
 Show all work.
\item
Let~\(f\) be Lebesgue integrable on \([0,1]\). Suppose that \(\int_0^x f\,dm=0\), for every \(x\in[0,1]\). What can you conclude about~\(f\)? Prove your answer?
\end{examproblems}
\end{document}
